% a03_Uplus: proof of prop:k02:Uplus, the two small-field lines at a non-negative scalar source
\section{Proof of \cref{prop:k02:Uplus}: the small-field lines at a non-negative scalar source}\label{app:a03_Uplus}

This appendix proves \cref{prop:k02:Uplus}: on King's cube $\CK$, at a real source $s=(u,h)$ with $h\ge0$, the lower line $\Zk(s)\ge e^{-C\eta^\sigma}\Ik{j}(s)$ and the small-field half of the upper line, $\int P\,d\mu_{k,s}\le e^{C\eta^\sigma}\Ik{j}(s)$, hold on both members of every diagonal pair, uniformly on compact source families. The large-field remainder at $s\ne0$ is not part of the statement; it is supplied where needed by \cref{lem:k02:HLF}.

The proof is the insertion of the scalar source into the small-field chain of the proof of \cref{prop:k01:H0rel}, that is, into Ba{\l}aban's inductive lower bound \cite[Sect.~3]{Ba1} and into the term of his upper bound \cite[(2.43)]{Ba2} in which no large-field block occurs at any scale, both run on $\CK$ with King's flow of \cref{ssec:k01:setting} and with the curvature source $u$ already present. The point is that the scalar source enters King's effective action through one vertex only, \cite[(3.21)]{K1}, whose coupling carries two powers of the running scale where the other vertices carry at most one and a half; every step of the chain is either blind to it, or a bound on the interaction families which that margin leaves unchanged, or a piece of bookkeeping that is written out below. The notation is that of \cref{ssec:k01:setting,ssec:k02:setting}.

\subsection{Scales, units, and the source vertex}\label{ssec:a03:setting}

Fix the couplings, $J\ge J_0$ and a spacing $\eps_k=5^{-k}$. At step $i$ of King's flow ($0\le i\le j$) the block side is $\ell_i\defeq5^{i}\eps_k$; the blocks of side $\ell_i$ tiling $\XJ$ form the lattice $\Lambda^{(i)}$, $|\Lambda^{(i)}|=\ell_i^{-3}|\CJ|$; and $p_i\defeq p(\ell_i)=b_0(1+\ln\ell_i^{-1})^{p}$ is Ba{\l}aban's threshold at that scale, with $p=6$ as in \cite[(3.54)]{Ba1}. In \emph{unit-$i$ units} lengths are measured in $\ell_i$, the fine lattice has spacing $5^{-i}$ and $\sum_{x\in B}5^{-3i}=1$ over the fine sites of a block $B$. At the stopping scale $i=j$ one has $\ell_j=\eta$; on a diagonal pair $\eta=\eta_k$.

The sources are supported at $\ell^\infty$-distance $>1$ from $\R^3\setminus\CJ$, so in unit-$i$ units every source vertex lies at distance $>1/\ell_i$ from the Dirichlet faces. We write $\mathsf S(s)=\sum_i|x_i|\max_{\mu\nu}\norm{f_{i,\mu\nu}}_{C^{2,\alpha}}+\norm h_\infty$ as in \cref{ssec:k02:inputs} and fix $S_0$ with $\mathsf S(s)\le S_0$.

\paragraph{The source vertex.} In King's generating function \cite[(3.17)]{K1} the scalar source appears as the vertex \cite[(3.21)]{K1}, which at step $i$ reads, in unit-$i$ units,
\begin{equation}\label{eq:a03:V21}
\theta\,\ell_i^{2}\sum_{x}5^{-3i}\,h(x)\sum_{a=1}^{N}\Big[\big(\phi'_a(x)+\phi^{(i)}_a(x)\big)^2-C^{(i)}(x,x)\Big],
\end{equation}
with $\theta$ the amplitude parameter of the vertex ($\theta=1$ in the action), $\phi'$ the fluctuation field, $\phi^{(i)}$ the background (external) scalar field of the step, and $C^{(i)}(x,x)$ King's Wick constant, the diagonal of the free scalar covariance at the running scale ($(-\Delta^{\eta,\Omega}+m^2(5^i\eps_k)^2)^{-1}(x,x)$ in King's notation; at $i=0$ it is the free Wick constant $\sW_{\eps_k}$ of \cref{def:s02:sources}). Its \emph{coupling power} is $2$: the other vertices of King's list \cite[(3.18)--(3.26)]{K1} carry $\lambda\ell_i$ (the quartic), $(e\ell_i^{1/2})^m$, $m\ge1$ (the current vertices), $e^2\ell_i$ (the seagull), $\tau\ell_i^{3/2}$ (the curvature source \cite[(3.20)]{K1}) and $O(\ell_i)$ (the mass counterterm, by \cite[(1.29)]{Ba3}). The vertex \eqref{eq:a03:V21} carries no photon field, is gauge invariant and charge neutral, and has at most two scalar legs.

\subsection{Where the sources sit}\label{ssec:a03:A}

\begin{lemma}[where the sources sit]\label{lem:a03:A}
Let $S^{(i)}_k(\cdot;s)$ be King's effective action after $i$ steps on $\CK$ at the source $s=(u,h)$, and $P^{(i)}_k(\cdot;s)$ its interaction part \cite[(3.17)]{K1}.
\begin{enumerate}[label=\textup{(\alph*)},nosep]
\item For every $i$ and $s$, $S^{(i)}_k(\cdot;s)-S^{(i)}_k(\cdot;0)=P^{(i)}_k(\cdot;s)-P^{(i)}_k(\cdot;0)$: the Gaussian parts $\frac12\ip{A_i}{\Delta^{(i)}A_i}$, $\frac12\ip{\phi_i}{\Delta^{(i)}(A^{(i)})\phi_i}$, the normalisation constants $\ln Z_i$, $\ln Z_i(A^{(i)})$ and the vacuum counterterms $E_0$, $E_1$ are source-free.
\item The indicator $\chi_i$ (King's \cite[(3.2)]{K1} times the face factor \cite[(5.41)]{K2}), every test of Ba{\l}aban's decomposition of unity at every scale \cite[(2.2)--(2.10)]{Ba2}, and the small-field classes inserted at the intermediate scales of the lower bound \cite[Sect.~3]{Ba1} are source-free functions of the block and fluctuation fields. Hence the product $P$ of the small-field factors, the admissible histories of large-field blocks and the region geometry are source-free.
\item Every renormalisation-group operator (the averaging operators, the block covariances $\Delta^{(i)}$, $C^{(i)}$, the propagators $G_i$, the backgrounds, the sub-domain operators and the difference lines between propagator domains), every object of the photon sector, every positivity statement and every Gaussian integration of the chain is source-free.
\item At step $0$ the scalar source enters the bare interaction as $-\eps_k^2h(x)\big[|\phi(x)|^2-N\,c_W(x)\big]$ in the exponent, in unit-$\eps_k$ units ($\eps_k^3h|\phi_{\mathrm{phys}}|^2=\eps_k^2h|\phi_{\mathrm{unit}}|^2$ with $\phi_{\mathrm{phys}}=\eps_k^{-1/2}\phi_{\mathrm{unit}}$; $c_W\defeq\eps_k\sW_{\eps_k}$ is a unit-lattice Wick constant, $0<c_W\le C_W$).
\item Consequently a step of the chain is touched by the scalar source only if it bounds, expands or removes $P^{(i)}_k$ or the bare interaction, or assembles such bounds. Every other step is $h$-blind.
\end{enumerate}
\end{lemma}

\begin{proof}
(a) King's \cite[(3.14)]{K1} reads $S^{(k),1}(T^{(k)},A_k,\phi_k,g,h)=\frac12\ip{A_k}{\Delta^{(k)}A_k}-\ln Z_k+\frac12\ip{\phi_k}{\Delta^{(k)}(A^{(k)})\phi_k}-\ln Z_k(A^{(k)})+E_0+P^{(k),1}(T^{(k)},A_k,\phi_k,g,h)$, with $E_0$ given by \cite[(3.15)]{K1}, a free Gaussian integral; the sources $g,h$ occur in $P^{(k),1}$ alone. By \cite[(3.17)]{K1}, $P^{(k),1}$ is a sum of derivatives with respect to the parameters $e',\lambda',\tau,\theta$ of $E'_k=-\ln\int\prod_jd\mu_{C(j),L^j\eta}(A'_j)\,d\mu_{G_k(A^{(k)})}(\phi')\,e^{V^{(0),\eta}}$, whose fluctuation measures are source-free Gaussians, and in the vertex list \cite[(3.18)--(3.26)]{K1} of $V^{(0),\eta}$ the curvature source occurs only in (3.20) and the scalar source only in (3.21). The same holds for $E_1$, King's \cite[(3.35)]{K1} on the cube, which is built from the source-free part of the same expansion. (b) \cite[(3.2)]{K1} and \cite[(5.41)]{K2} are inequalities on $(A_i,\phi_i)$; Ba{\l}aban's tests \cite[(2.2)--(2.10)]{Ba2} and the classes of \cite[Sect.~3]{Ba1} (King's \cite[(A.1)]{K1}) are inequalities on the fields. (c) The operators are built from the Gaussian forms and the averaging operators, which carry no source; the photon sector does not see $h$ (\eqref{eq:a03:V21} has no photon field); the positivity statements and Gaussian integrations are statements about the source-free kernels. (d) \cite[(3.21)]{K1} at $k=0$ is the bare source with King's Wick constant; Ba{\l}aban's bare interaction \cite[p.~559]{Ba2} is $\lambda(\eps)|\phi|^4+\frac12\dm\,\eps^2|\phi|^2$, to which the source adds the displayed term. (e) is the summary of (a)--(d).
\end{proof}

\subsection{Power counting with the source vertex}\label{ssec:a03:PC}

In the power counting of \cite[Sect.~3]{K1} and \cite[Sects.~1--3]{Ba3} a vertex $v$ whose coupling carries $\ell_i^{u_v}$ in unit-$i$ units contributes $u_v$, and a connected graph $\gamma$ with $E_\gamma$ external legs (scalar or photon) has degree
\begin{equation}\label{eq:a03:degree}
D(\gamma)\ \defeq\ \sum_{v\in\gamma}u_v-3+\tfrac12E_\gamma ;
\end{equation}
the subgraphs that are renormalised (the divergent ones) are those with $D(\gamma)\le0$. The quartic has $u=1$, the current vertex of order $m$ has $u=m/2$, the seagull $u=1$, the mass counterterm $u\ge1$, the curvature vertex $u=3/2$, the remainder vertices of order $\bar n+1=14$ have $u\ge7$, and the source vertex \eqref{eq:a03:V21} has $u=2$.

\begin{lemma}[the divergent graphs containing a source vertex]\label{lem:a03:PC}
Let $\gamma$ be a connected graph with at least one line.
\begin{enumerate}[label=\textup{(\roman*)},nosep]
\item If $\gamma$ contains exactly one vertex \eqref{eq:a03:V21} and no curvature vertex, then $D(\gamma)\le0$ only in the following cases: \textup{(T)} $E_\gamma=0$, the graph is a field-free $h$-linear vacuum graph whose other vertices have total coupling power $\le1$, that is, of order $(\alpha,\beta)\in\{(0,0),(0,1),(2,0)\}$ in $(e^2,\lambda)$, with the mass counterterm counted as its partner; \textup{(F)} $\gamma=\gamma_F$, the graph with the source vertex at $x$, one current vertex of order $m=1$ at $y$, the two scalar lines joining $x$ and $y$, and the photon leg of $y$ external; $D(\gamma_F)=0$, and its amplitude vanishes identically.
\item If $\gamma$ contains two or more vertices \eqref{eq:a03:V21}, or one of them and a curvature vertex, then $D(\gamma)\ge\frac12$.
\end{enumerate}
Consequently no subgraph containing a source vertex and an external leg is renormalised, and the renormalisation of the $h$-free divergent subgraphs inside a family with a source vertex is the one at $h=0$.
\end{lemma}

\begin{proof}
(i) With one source vertex, $D(\gamma)=\sum_{v\ne\text{source}}u_v-1+\frac12E_\gamma$, so $D\le0$ forces $\sum_{v\ne\text{source}}u_v+\frac12E_\gamma\le1$. If $E_\gamma=0$ the other vertices have total coupling power $\le1$: none (the Wick tadpole, order $(0,0)$), one quartic or one mass counterterm (order $(0,1)$), or two current vertices of order $1$, or one seagull, or one current vertex of order $2$ (order $(2,0)$); a single current vertex of order $1$ is excluded at $E_\gamma=0$, its photon leg having no partner. If $E_\gamma=1$ the other vertices have total power $\le\frac12$, so there is exactly one current vertex of order $1$; the external leg is its photon leg (a lone scalar leg cannot be external when the source vertex's two legs are both to be contracted and the current vertex's two scalar legs are the only other ones), both scalar legs of each vertex are contracted into lines, and the graph is $\gamma_F$. If $E_\gamma=2$ there is no other vertex and no line, which is excluded. If $E_\gamma\ge3$ then $D\ge\frac12$. The amplitude of $\gamma_F$ is the scalar loop with one insertion of $|\phi|^2$ and one insertion of the current $e\,A\cdot(\phi,q\,\partial\phi)$, whose component trace is $\Tr q$ with $q$ the antisymmetric charge generator of the $U(1)$ action on $\R^N$; $\Tr q=0$, so $\gamma_F$ vanishes. (ii) Two source vertices give $\sum_vu_v\ge4$, hence $D\ge1$; a source vertex and a curvature vertex give $\sum_vu_v\ge\frac72$, hence $D\ge\frac12$. The consequence: by (i), (ii), the only divergent subgraphs with a source vertex are vacuum graphs ($E=0$), which are not Taylor-subtracted at external legs, and $\gamma_F$, which is zero; the divergent subgraphs with external legs are therefore source-free, and they are renormalised as at $h=0$.
\end{proof}

The three field-free $h$-linear families of case (T) are exactly the members excluded from the class $\mathcal N_h$ of \cref{hyp:k02:Kh}. In the small-field chain they are never bounded: they are field-independent constants of the effective action, carried exactly (\cref{ssec:a03:SF}, step B7).

\subsection{The margin}\label{ssec:a03:D}

The graph bound used throughout the chain is Ba{\l}aban's bound in product form \cite[Prop.~1, (1.33)]{Ba3}, King's \cite[Thm.~3.5, (3.38)]{K1} on the cube: in unit-$i$ units, on Ba{\l}aban's small-field set, every renormalised connected family of $P^{(i)}_k$ with at least one external field has a coefficient bound, per block, $O(1)\,\ell_i^{\rho}\,p_i^{n_1}e^{-\delta_0d}$, where $d$ is the tree length of the vertex positions and the power $\rho$ is the sum over the vertices of the coupling power $u_v$, less $\frac14$ for every external scalar leg: on Ba{\l}aban's intermediate field set the external scalar field is bounded by $|\psi(y)|\le2\lambda(\eps)^{-1/4}p(\eps)$ \cite[p.~559]{Ba2}, that is by $C\,p_i\,\ell_i^{-1/4}$ in unit-$i$ units, while lines and photon legs cost no power. Writing $\epsilon_v\defeq u_v-\frac14\,(\text{number of external scalar legs of }v)$, the power is $\rho=\sum_v\epsilon_v$, and $\epsilon_v\ge0$ for every vertex of the list \cite[(1.6)--(1.15)]{Ba3}. For the families with a source vertex this bound is the cube form of the (3.38) bound of \cref{hyp:k02:Kh}; the vertex \eqref{eq:a03:V21} is admitted into the list as an ordinary local two-leg vertex with coefficient bounded by $\norm h_\infty\le S_0$, since by \cref{lem:a03:PC} it creates no new divergent subgraph with external legs, so that the renormalisation structure of the families, and with it the proof of the product-form bound, is the one at $h=0$.

\begin{lemma}[the margin of the source families]\label{lem:a03:D}
In unit-$i$ units, on Ba{\l}aban's step-$i$ intermediate field set, every renormalised connected family of $P^{(i)}_k(\cdot;s)$ with at least one source vertex \eqref{eq:a03:V21} and at least one external field (scalar, block or photon) has a coefficient bound, per block,
\[
C(S_0)\,\ell_i^{3/2}\,p_i^{N'}\,(1+\log_5\ell_i^{-1})^{p_0}\ \le\ C'(S_0)\,\ell_i^{5/4},
\]
with $C'(S_0)$ independent of $i$, $k$ and of $s$ with $\mathsf S(s)\le S_0$. The same holds for the families with a curvature vertex and no source vertex.
\end{lemma}

\begin{proof}
The source vertex has coupling $\theta\ell_i^2h$ and at most two scalar legs, so $\epsilon_v\ge2-\frac12=\frac32$; the curvature vertex \cite[(3.20)]{K1} has $\epsilon_v=\frac32$ (no scalar leg); every other vertex has $\epsilon_v\ge0$. Hence $\sum_v\epsilon_v\ge\frac32$ for every family containing a source or a curvature vertex. By \cref{lem:a03:PC} no Taylor subtraction sits at a subgraph containing the source vertex, so the product-form bound applies with \eqref{eq:a03:V21} as a local vertex of sup norm $\le S_0$; the intermediate set enlarges the leg bounds of King's set \cite[(3.2)]{K1} by a fixed factor, which is absorbed into $C(S_0)$; and $p_i^{N'}(1+\log_5\ell_i^{-1})^{p_0}\le C\ell_i^{-1/4}$. The constants depend on $s$ only through $\norm h_\infty$ and $\max\norm{f_{i,\mu\nu}}_{C^{2,\alpha}}$, hence through $S_0$. Near the Dirichlet faces: the source vertices lie at unit-$i$ distance $>1/\ell_i$ from $\R^3\setminus\CJ$, so a family with a source vertex and a vertex within the face zone has tree length $d>1/\ell_i$ and carries the additional factor $e^{-\delta_0d}\le e^{-\delta_0/(2\ell_i)}\le\ell_i^{10}$ beyond the face bounds of the proof of \cref{prop:k01:H0rel}; it enters every estimate through this tail only.
\end{proof}

The consequence used below is that the uniform exponent $\beta_0=\frac14-\beta$ of the coefficient bounds \cite[Prop.~3.2]{Ba1} and the summability condition $(\bar n+1)\beta_0>3$ of the cumulant step hold at $s$ unchanged, since $\frac54\ge\beta_0$; and that the remainder terms of order $>\bar n$ and the terms of total power $>3$ in $\ell_i$ \cite[Prop.~3.1]{Ba1} keep their bound $O_{S_0}(\ell_i^{\kappa'})|\Lambda^{(i)}|$ with the same $\kappa'>3$.

\subsection{The two small-field lines}\label{ssec:a03:SF}

\begin{theorem}[the small-field lines at $h\ge0$]\label{thm:a03:SF}
Fix $J\ge J_0$, $\alpha\in(0,\frac12)$ and $S_0>0$. There are $k_0$, $C(S_0)$ and $\sigma=\kappa'-3>0$, depending on the couplings, $J$, $\alpha$, $S_0$ and on nothing else, such that for every real source $s=(u,h)$ with $h\ge0$ and $\mathsf S(s)\le S_0$, every $k\ge k_0$ and both members $(k',j')\in\{(k,j(k)),(k+1,j(k)+1)\}$ of the diagonal pair, with $b\defeq C(S_0)|\CJ|\,\eta_k^{\sigma}$,
\[
Z_{k'}(s)\ \ge\ e^{-b}\,I^{(j')}_{k'}(s)\qquad\text{and}\qquad\int P\,d\mu_{k',s}\ \le\ e^{b}\,I^{(j')}_{k'}(s).
\]
\end{theorem}
\noindent Since $C(S_0)|\CJ|$ depends on $J$ and $S_0$ only, this is \eqref{eq:k02:SF}, and \cref{prop:k02:Uplus} follows.

\begin{proof}
We fix one member, write $k$ for its spacing index and $j$ for its number of steps, and run the small-field chain of the proof of \cref{prop:k01:H0rel} at the source $s$ instead of $(u,0)$. By \cref{lem:a03:A}(e) only the steps that bound, expand or remove the interaction, or assemble such bounds, see the change; we go through them.

\emph{The lower line} is Ba{\l}aban's inductive lower bound \cite[Sect.~3]{Ba1}, in the following seven steps.

\textit{B1 (restriction to small fields, \cite[(3.6)]{Ba1}).} The integrand $e^{-S}e^{F(u)+\Wick{h}}$ of $\Zk(s)$ is positive, the sources being real. Inserting the source-free indicators (\cref{lem:a03:A}(b)) of the small-field classes at the scales $i<j$ and of $\chi_j$ at the stopping scale gives a lower bound. Nothing is paid for in a lower bound.

\textit{B2 (the inductive hypothesis, \cite[(3.26)--(3.32)]{Ba1}).} After $i$ steps, $\Zk(s)\ge e^{-a_i}\int\chi_{(i)}e^{-S^{(i)}_k(\cdot;s)}$ with $a_i\defeq\sum_{i'<i}C(S_0)\ell_{i'}^{\kappa'}|\Lambda^{(i')}|$, where $S^{(i)}_k(\cdot;s)$ is King's effective action whose vertex list contains \cite[(3.20)--(3.21)]{K1} from the start, and $\chi_{(i)}$ is the inserted class indicator. The identities for the normalisation constants and for the representation of the fields after $i$ steps are source-free (\cref{lem:a03:A}(c)).

\textit{B3 (the fluctuation step, \cite[(3.37)--(3.49), Props.~3.1--3.2]{Ba1}).} The Gaussian fluctuation integral is exact and source-free. The interaction, now with the source families, is expanded to order $\bar n=13$ with Cauchy estimates in the complexified parameters $(e',\lambda',\tau,\theta)$ of \cite[(3.17)]{K1}; $\theta$ is the amplitude of \eqref{eq:a03:V21}, and on $|\theta|\le2$ the coupling of \eqref{eq:a03:V21} is at most $2\ell_i^2S_0$, so \cref{lem:a03:D} applies with $S_0$ replaced by $2S_0$. The remainder terms and the terms of total power $>3$ are bounded by $O_{S_0}(\ell_i^{\kappa'})|\Lambda^{(i)}|$, $\kappa'>3$, and the coefficients of the polynomial part by $O_{S_0}(\ell_i^{\beta_0})$ per block, as at $h=0$: by \cref{lem:a03:D} every family with a source vertex and an external field has coefficient $O(\ell_i^{5/4})$, within both bounds. The one-loop determinants are source-free.

\textit{B4 (cutoff removal and the consistency conditions, \cite[(3.50)--(3.55)]{Ba1}).} The cumulant bound used at this step in the proof of \cref{prop:k01:H0rel} is blind to the weights and to the vertex list: it needs only that the interaction be a real polynomial in the fluctuation variables without constant term, with coefficients bounded as in B3. At the source $s$ the fluctuation-independent parts of the expanded interaction, namely the background part $\theta\ell_i^2\sum5^{-3i}h|\phi^{(i)}|^2$ of \eqref{eq:a03:V21}, the Wick member, and every field-free family, are pulled out of the exponent exactly, into $S^{(i+1)}_k(\cdot;s)$; what remains is a polynomial in the fluctuation without constant term, with coefficients as in B3. The compatibility inequality \cite[(3.53)]{Ba1} between the classes of consecutive scales is a statement on fields (\cref{lem:a03:A}(b)), and the conditions \cite[(3.54)]{Ba1} on $p$ are unchanged.

\textit{B5 (the classes).} The background classes inserted at the intermediate scales are source-free.

\textit{B6 (the final step, \cite[(3.66)--(3.69)]{Ba1}).} It gives the lower bound $\Zk(0)\ge e^{-E_-|\CJ|}$ on the partition function at zero source; it is used in the proof of \cref{prop:k01:H0rel} at $s=0$ only and is not needed here.

\textit{B7 (the sources).} The curvature source is King's plaquette source \cite[(2.7)]{K2}. The source vertex \eqref{eq:a03:V21} is switched on. Its field-independent part, $-\theta\ell_i^2\sum_x5^{-3i}h(x)\,N\,C^{(i)}(x,x)$ together with every field-free family it generates, is a $c$-number of $S^{(i)}_k(\cdot;s)$ at every step. A field- and fluctuation-independent constant in the action multiplies the integral and every intermediate integrand by the same factor; so in the small-field half it is common to $\Zk(s)$ and to $\Ik{j}(s)$, with King's Wick constant inside $S^{(j)}_k(\cdot;s)$ on both sides. No bound on these constants is needed, and none is used.

\textit{Assembly (\cite[Sect.~3, end]{Ba1}).} After $j$ steps, $\Zk(s)\ge e^{-a}\int\chi_je^{-S^{(j)}_k(\cdot;s)}=e^{-a}\Ik{j}(s)$ with
\[
a=\sum_{i<j}C(S_0)\,\ell_i^{\kappa'}|\Lambda^{(i)}|=C(S_0)|\CJ|\sum_{i<j}\ell_i^{\kappa'-3}\ \le\ C(S_0)|\CJ|\,\eta^{\sigma},\qquad\sigma=\kappa'-3>0,
\]
the sum being geometric and dominated by $i=j-1$, $\ell_{j-1}=\eta/5$; the constant does not depend on $k$ or $j$. The restriction to $\chi_j$ at the last step is a positivity step inside the class, so the thresholds of $\chi_j$ never enter a constant.

\emph{The upper line's small-field part.} In Ba{\l}aban's inductive formula \cite[(2.43)]{Ba2} at the scale $j$, with the decomposition of unity of \cite[(2.2)--(2.10)]{Ba2} inserted at every scale $i\le j$ (source-free, \cref{lem:a03:A}(b)), $\Zk(s)$ is a sum over histories of large-field blocks, and $\int P\,d\mu_{k,s}$ is the term of the history with no large-field block at any scale $i\le j$. Along this history the steps of the upper bound (the restriction to the small-field set, the translation to the background, the expansion of the interaction to order $\bar n$ with its cumulants, the bounds on the background and fluctuation fields, the $Z$-factors and the Gaussian integrations, \cite[Sect.~2]{Ba2}) are carried out with the source families present and nothing is removed. Of these steps only the expansion of the interaction sees the source; it gives, by \cref{lem:a03:D} exactly as in B3, the same remainder bound $O_{S_0}(\ell_i^{\kappa'})|\Lambda^{(i)}|$ at every scale $i<j$, and the field-free constants are carried exactly as in B7. The replacement of Ba{\l}aban's cut photon background by the full one in the final integrand is an operation in the photon sector, source-free, done as in the proof of \cref{prop:k01:H0rel}. Hence
\[
\int P\,d\mu_{k,s}\ \le\ \int\chi_j\,e^{-S^{(j)}_k(\cdot;s)+\sum_{i<j}C(S_0)\ell_i^{\kappa'}|\Lambda^{(i)}|}\ =\ e^{a}\,\Ik{j}(s),
\]
with the same $a$ as in the lower line; the last step uses that $\chi_j$ is the indicator of $\Ik{j}$ and that the integrand is positive, the exponent being real.

\emph{Both members.} Both members of a diagonal pair have the stopping scale $\eta_k$, the same unit lattice and the same indicator $\chi_{j(k)}$ (\cref{ssec:k01:setting}); every constant above depends on $(k,j)$ only through $\eta$, so one $b=C(S_0)|\CJ|\eta_k^\sigma$ serves both members, and $\Ik{j}(s)$ is the same kind of integral on both.

\emph{Uniformity on compact families.} The source enters every constant only through $\norm h_\infty$, by way of the coupling of \eqref{eq:a03:V21} (no Taylor subtraction sits at a source subgraph, \cref{lem:a03:PC}, so no derivative of $h$ is ever taken), and through $\max_{\mu\nu}\norm{f_{i,\mu\nu}}_{C^{2,\alpha}}$ as in the proof of \cref{prop:k01:H0rel}. The constants are non-decreasing functions of $S_0$; so the statement is uniform on $\{\mathsf S(s)\le S_0\}$. The sign $h\ge0$ is not used in this proof; it is the hypothesis under which \cref{lem:k02:HLF} supplies the large-field remainder, and it is kept in the statement because the proposition is consumed only there.

\emph{Every point of King's range.} Each step holds at every point $(e,\lambda,m,\lmu)$ of King's range with point-dependent constants, in particular at the mass $M_2$ of \cref{ssec:k02:setting} with the same $e,\lambda,\lmu$.
\end{proof}

\subsection{The insertion, step by step}\label{ssec:a03:census}

For completeness we list how each step of the small-field chain of the proof of \cref{prop:k01:H0rel} is affected by the scalar source; the marks are: \textbf{blind}, the step is source-free by \cref{lem:a03:A}; \textbf{family}, the step is a bound on the families of the interaction, read at $s$ through \cref{hyp:k02:Kh} and \cref{lem:a03:D}; \textbf{rewritten}, the step's bookkeeping changes at $s$ and is written out in \cref{ssec:a03:SF}.

\begin{center}\small
\begin{tabular}{p{0.46\textwidth}lp{0.38\textwidth}}
\hline
step & mark & reason \\
\hline
the operator layer: block covariances, propagators and their decay, sub-domain operators, difference lines, Green's function regularity \cite{Ba4} & blind & $h$ is a vertex; it enters no Gaussian form \\
the graph bound in product form \cite[Prop.~1]{Ba3}, King's \cite[(3.38)]{K1}, on the cube & family & source families by \cref{hyp:k02:Kh}; the margin by \cref{lem:a03:D}; the face zone through tails only \\
the Taylor step at a divergent two-leg subgraph & family & no subtraction at a source subgraph (\cref{lem:a03:PC}; $\gamma_F$ vanishes) \\
the Ward identities for the renormalised classes & family & \eqref{eq:a03:V21} is gauge invariant; the identities are unchanged \\
the bounds on propagator sub-domains and the domain-change estimate \cite[Prop.~2.6]{Ba2}, \cref{hyp:k01:BALint} & family & the same reading with \eqref{eq:a03:V21} in the list \\
\cite[Props.~3.1--3.2]{Ba1}: remainder terms and coefficient bounds & family & at $s$ with \cref{lem:a03:D}; exponents unchanged \\
the uniform exponent $\beta_0$ & family & \cref{lem:a03:D}: $\epsilon_v\ge\frac32$ at the source vertex \\
the small-field tests, the classes, their compatibility, the external-field properties & blind & functions of the fields \\
B1 positivity & rewritten & real sources \\
B2 inductive hypothesis & rewritten & $S^{(i)}_k(\cdot;s)$ \\
B3 fluctuation step & family & Cauchy estimates also in $\theta$ \\
B4 cutoff removal & family & constants pulled out exactly \\
B5 classes & blind & \\
B6 final step & blind & consumed at $s=0$ only \\
B7 sources & rewritten & the source vertex switched on \\
lower-line assembly & rewritten & the lower line at $s$ \\
positivity lemmas, the decomposition of unity, the region geometry & blind & \\
the restriction to small fields and the translation to the background in the upper bound \cite[(2.11)--(2.14), (2.43)]{Ba2} & rewritten & the bare source re-expressed after the translation, exactly \\
the expansion of the interaction in the upper bound & rewritten & with the source families; nothing removed on the small-field history \\
backgrounds, fluctuation bounds, $Z$-factors, last-step masses, positivity, diamagnetic integration & blind & \\
the per-block costs of the final scale and the final-scale bound & blind & consumed in the large-field remainder only \\
the upper-line assembly & rewritten & the small-field part at $s$ \\
\hline
\end{tabular}
\end{center}

Also outside the enumerated steps: the basic positivity lemma, the composition identities of the flow and the Schur complements are blind; the definition of the effective action carries the vertex list with \eqref{eq:a03:V21} switched on; the $c$-number lemma (a field- and source-independent constant multiplies every integrand by the same factor) is used at B7 with King's Wick constant in place of a scheme constant; the face lemmas of the cube (the Dirichlet mismatch of the counterterm, the face terms of the photon sector) are blind, carrying no source vertex.
